\documentclass[11pt]{amsart}
\usepackage[T1]{fontenc}
\usepackage[margin=1.12in]{geometry}
\usepackage{amsmath,amssymb,amsthm,mathtools,microtype}
\usepackage{xcolor}
\definecolor{linkblue}{RGB}{25,55,110}
\usepackage[colorlinks=true,linkcolor=linkblue,citecolor=linkblue,urlcolor=linkblue]{hyperref}
\hypersetup{pdftitle={An obstruction to uniform residual-twist compression},pdfauthor={Caleb Barnett}}
\newtheorem{proposition}{Proposition}[section]
\newtheorem{lemma}[proposition]{Lemma}
\theoremstyle{remark}
\newtheorem{remark}[proposition]{Remark}
\newcommand{\Z}{\mathbb Z}
\newcommand{\degp}{\deg_p}
\title[Uniform residual-twist compression]{An obstruction to the uniform residual-twist compression in Carmi--Cohen}
\author{Caleb Barnett}
\urladdr{https://barnettcaleb1.github.io}
\date{First complete draft for internal review: 3 November 2026}
\subjclass[2020]{57K10, 57K14}
\keywords{Jones polynomial, finite-type invariant, twist knot}

\begin{document}
\begin{abstract}
We give an exact obstruction to the length-independent residual-twist
bound used in Carmi and Cohen, \emph{On trivial Jones--Vassiliev
polynomials}, arXiv:2606.22410v1. An ordinary antiparallel twist family
based at a minimal trefoil diagram has Jones--Vassiliev degree \(2L+5\)
at odd twist length \(L\), with leading coefficient \(-1\).
The asserted geometric package, followed by Theorem 5.6 and Lemma 5.8
of that preprint, would bound these degrees by \(22\) independently
of \(L\). The permitted value \(L=23\) instead has \(b_{51}=-1\).
Thus the universal geometric package, the signed-state compression
conclusion, and the ordinary-residual identification cannot all hold.
This note does not refute the paper's main
statement or the Jones unknot conjecture.
\end{abstract}
\maketitle

\section{Scope and conventions}

This note concerns the version of Carmi--Cohen \cite{CC} submitted on
21 June 2026. All its result numbers below refer to that version.
The issue is a uniform-control assertion in its proof, not a proposed
counterexample to Jones unknot detection. The knots used below are
ordinary twist knots, and their Jones polynomials are explicitly
nontrivial.

We use the preprint's variables and normalization:
\[
 x=t^{1/2},\qquad p=x-x^{-1},\qquad V_U=1,
\]
with the oriented skein relation
\begin{equation}\label{eq:skein}
 x^{-2}V_+-x^2V_-=pV_0.
\end{equation}
There is a ring isomorphism
\begin{equation}\label{eq:ring}
 \Z[x,x^{-1}]\cong\Z[p,x]/(x^2-px-1).
\end{equation}
Indeed \(x^{-1}=x-p\) in the quotient, and \(p=x-x^{-1}\)
recovers the inverse map. The quotient is free over \(\Z[p]\)
with basis \(1,x\), since its defining polynomial is monic of
degree two in \(x\). Every Jones polynomial therefore has a
unique expression
\begin{equation}\label{eq:jvp}
 V_K=a_K(p)+x b_K(p),\qquad
 a_K=\sum_r a_r(K)p^r,\quad b_K=\sum_r b_r(K)p^r.
\end{equation}
The degree discussed here is
\(\degp V_K=\max(\deg a_K,\deg b_K)\).
Write \(f_r(K)=(a_r(K),b_r(K))\). In the source,
\(G_r(\mathcal Q)\) means that \(f_r\) vanishes on every
state in the indicated cube \(\mathcal Q\).

The source's main conclusion asserts a nonzero Jones finite-type
coefficient, at order at most three times crossing number, for
every nontrivial knot. The calculation below neither proves nor
disproves that assertion. It contradicts a different assertion
used to obtain it: a bound on \emph{all} sufficiently high
residual-twist coefficient layers, uniform in the twist length
and depending only on the original host.

\section{The ordinary twist family}\label{sec:family}

Take a three-crossing trefoil in its twist-knot diagram: a
two-crossing clasp and one marked antiparallel seed crossing.
Replace the seed by \(L\) crossings of the same sign, where
\(L\ge1\) is odd, and call the result \(K_L\).
The parity retains the endpoint pairing, so every \(K_L\)
is a knot. Choose the mirror convention for which
\begin{equation}\label{eq:seed}
 V_1=x^2+x^6-x^8.
\end{equation}
The family can also be specified as the appropriate mirror of
the two-bridge diagram with continued fraction \([L,2]\),
namely rational fraction \((2L+1)/2\);
equation~\eqref{eq:seed} fixes the chirality.
The changing \(L\)-box is the marked seed, while the
two-crossing clasp is fixed.

Switching the seed in \(K_1\) gives the unknot. Let \(H\)
be the Jones polynomial of its oriented smoothing.
The seed skein relation determines it without an independent
sign choice:
\[
 pH=x^{-2}(x^2+x^6-x^8)-x^2=1-x^2+x^4-x^6.
\]
Consequently
\begin{equation}\label{eq:smoothing}
 H=-x-x^5.
\end{equation}
Changing one crossing in the positive antiparallel \(L\)-box
allows a Reidemeister-II cancellation of two crossings,
leaving \(K_{L-2}\). Oriented smoothing makes the other
crossings in that box removable curls, and leaves the same
oriented smoothing link as in \eqref{eq:smoothing}.
The normalized Jones polynomial is invariant under their
removal. It follows that for odd \(L\ge3\)
\begin{equation}\label{eq:rec}
 V_L=x^4V_{L-2}+x^2-x^4+x^6-x^8.
\end{equation}
The \(L=1\) seed relation has the same form with the
switched state \(V_{-1}=1\).

\begin{proposition}\label{prop:degree}
For every positive odd \(L\),
\[
 \max\deg_x V_L=2L+6,\qquad
 \degp V_L=2L+5,\qquad b_{2L+5}(K_L)=-1.
\]
\end{proposition}
\begin{proof}
The initial leading term is \(-x^8\).
At every subsequent step of \eqref{eq:rec},
multiplication by \(x^4\) raises its exponent by four;
the inhomogeneous summand has degree only eight
and cannot cancel it. Thus the unique highest
monomial is \(-x^{2L+6}\).
All exponents in \(V_L\) are positive.

For \(n\ge1\), write \(x^n=A_n(p)+xB_n(p)\).
Multiplication by \(x\) in \eqref{eq:ring} gives
\[
 (A_{n+1},B_{n+1})=(B_n,A_n+pB_n).
\]
Starting with \(A_1=0,B_1=1\), it follows that
\(B_n\) is monic of degree \(n-1\), and \(A_n\)
has degree at most \(n-2\).
Every smaller positive power of \(x\) therefore
has \(B\)-degree strictly less than \(2L+5\).
The leading coefficient from \(-x^{2L+6}\)
survives in \(b_{K_L}\), and equals \(-1\).
This proves all three assertions.
\end{proof}

For example \(K_3=5_2\), with the convention
\[
 V_3=x^2-x^4+2x^6-x^8+x^{10}-x^{12}.
\]
The seed and this value agree with the published
polynomials in \cite[pp.~74, 79]{GZ} under \(q=x^2\).
The appendix gives a closed formula, four initial
values, and exact checks independent of the
leading-term argument.

\section{The incompatible uniform bound}\label{sec:package}

The relevant implication in \cite{CC} uses three
numbered statements. Their hypotheses must be
retained when reading the contradiction.
\begin{enumerate}
\item Proposition~5.12 (statement on printed
pp.~39--40) asserts a local geometric verification
of the clasped-twist construction and its
conditions \((\mathrm C1)\)--\((\mathrm C4)\),
with fixed overhead \(\kappa_{\mathrm{mach}}=11\).
In particular condition \((\mathrm C4)\)
identifies the double-smoothed machine with the
ordinary residual antiparallel twist at the
host crossing, relative to the exterior.

\item Theorem~5.6 (printed pp.~30--31) assumes
\((\mathrm C1)\)--\((\mathrm C3)\) uniformly
over a fixed exterior host family \(\mathcal H\).
It asserts
\[
 G_r(Q_M(\mathcal H))\quad\text{for all }r>\widehat B(\mathcal H),
 \qquad \widehat B=B_{\mathrm{host}}+11.
\]
Its statement explicitly says:
\emph{The conclusion is independent of the
residual twist length \(L\).}

\item Lemma~5.8 (printed pp.~31--33) assumes
\((\mathrm C1)\)--\((\mathrm C4)\) and transfers
this vanishing to the entire ordinary residual
cube \(Q_\rho\):
\begin{equation}\label{eq:claimed}
 G_r(Q_\rho(\mathcal H))\quad
 \text{for every }r>\widehat B(\mathcal H)+2.
\end{equation}
This is not restricted to a two-state residual
face. In particular, it includes the all-positive
length-\(L\) state.
\end{enumerate}

These assertions precede the later anchored
descent in Theorem~5.10. Their stated hypotheses
do not require the host to have Jones polynomial
one, nor does Lemma~5.8 require
\(L>\widehat B+2\). The trefoil host is therefore
a legitimate test of the universal geometric
verification.

For a fixed host shadow with \(n\) crossings
and \(\ell\) components, the source uses
\(B_{\mathrm{host}}=3n+\ell-1\);
see its Corollary~3.7, Section~5.1.1
(printed p.~25), and Remark~5.7.
Take \(\mathcal H\) to be the crossing-flip family
of the fixed three-crossing trefoil shadow.
The pre-planting count includes the seed crossing.
The minimal trefoil host has \(n=3,\ell=1\),
so the combined assertions give
\begin{equation}\label{eq:cap22}
 \degp V_L\le3\cdot3+11+2=22
\end{equation}
for every allowed residual length.
Taking \(L=23\), Proposition~\ref{prop:degree}
instead gives
\begin{equation}\label{eq:contradiction}
 \degp V_{23}=51,\qquad b_{51}(K_{23})=-1.
\end{equation}
This length also satisfies \(L>\widehat B+2\),
so imposing that later length inequality would
not remove the example. Changing the overhead
eleven to another fixed constant cannot repair
the conclusion, because the actual degree
\(2L+5\) is unbounded.

\subsection*{Why this does not refute the conditional skein algebra}

The paired two-clasp identity used in Lemma~5.8 is
\begin{equation}\label{eq:paired}
 V_{-+}+V_{+-}=x^{-4}V_{++}+x^4V_{--}-p^2V_{00}.
\end{equation}
It follows by applying \eqref{eq:skein} in
the two marked coordinates. For example,
\begin{align*}
 V_{-+}&=x^{-4}V_{++}-p^2V_{00}-px^2V_{0-},\\
 V_{+-}&=x^4V_{--}+px^2V_{0-};
\end{align*}
adding proves \eqref{eq:paired}.
Multiplication by \(x^{4}\) or \(x^{-4}\)
raises the basis degree by at most four,
since \(x^{-1}=x-p\).
If the four signed machine states have degree
at most \(\widehat B\), the right side solved
for \(p^2V_{00}\) has degree at most
\(\widehat B+4\). Freeness of the basis \(1,x\)
then gives \(\degp V_{00}\le\widehat B+2\).
This conditional calculation is valid.

The contradiction therefore shows that the
universal geometric verification, the
signed-state compression conclusion, and
the asserted ordinary-residual identification
cannot all hold for the displayed family.
It is not an example independently satisfying
all of \((\mathrm C1)\)--\((\mathrm C3)\)
and violating a conditional theorem.
It also does not, by itself, isolate one
diagrammatic submove as the unique error.
The combined package is enough to expose
the gap in the proof as written.

\section{What remains open}

The main assertion of \cite{CC} is not
\eqref{eq:cap22}. It is a claim about the
existence of a nonzero low-order coefficient
for a nontrivial knot. This note has not
produced a knot violating that claim.
Every \(K_L\) has the explicit nonconstant
Jones polynomial computed above.
The number three in \eqref{eq:cap22}
is the crossing count of the \emph{fixed host},
not that of the length-\(L\) residual knot.
Confusing those two counts would misstate
the obstruction.
The ordinary diagram of \(K_L\) has \(L+2\)
crossings, so the general diagram bound from
Corollary~3.7 is \(3L+6\), consistent with
the actual degree \(2L+5\).
That general bound is not contradicted.

The consequence established here is that
the length-independent uniform-control
step needed by the submitted proof is false
as a universal package. A repair would
have to replace or restrict that step
and verify the new hypotheses. Neither
Jones unknot detection nor its negation
follows from this calculation.

\appendix
\section{Closed formula and exact verification}\label{app:verify}

Put \(m=(L-1)/2\) and \(B=x^2-x^4+x^6-x^8\).
Iteration of \eqref{eq:rec} gives the polynomial identity
\begin{equation}\label{eq:closed}
 V_{2m+1}=x^{4m}(x^2+x^6-x^8)
       +B\sum_{j=0}^{m-1}x^{4j},\qquad m\ge0,
\end{equation}
where the empty sum is zero.
Multiplication by \(1-x^4\) also gives an
equivalent expression without a sum.
The sum form makes clear that no
denominator specialization is involved.

The first four values are
\begin{align*}
 V_1={}&x^2+x^6-x^8,\\
 V_3={}&x^2-x^4+2x^6-x^8+x^{10}-x^{12},\\
 V_5={}&x^2-x^4+2x^6-2x^8+2x^{10}-x^{12}+x^{14}-x^{16},\\
 V_7={}&x^2-x^4+2x^6-2x^8+2x^{10}-2x^{12}
       +2x^{14}-x^{16}+x^{18}-x^{20}.
\end{align*}

For an independent coefficient extraction, define
\[
 F_0(p)=0,\qquad
 F_n(p)=\sum_{j=0}^{\lfloor(n-1)/2\rfloor}
 \binom{n-1-j}{j}p^{n-1-2j}\quad(n\ge1).
\]
Pascal's identity gives \(F_{n+1}=pF_n+F_{n-1}\),
and direct reduction in \eqref{eq:ring} gives
\begin{equation}\label{eq:binomial}
 x^n=F_{n-1}(p)+xF_n(p),\qquad n\ge1.
\end{equation}
Thus the two basis coefficients of
\eqref{eq:closed} can be obtained from
binomial coefficients, without iterating
the basis multiplication used in the
proof of Proposition~\ref{prop:degree}.
For instance,
\begin{align*}
 a_{K_1}(p)&=1-3p^2-4p^4-p^6,\\
 b_{K_1}(p)&=-6p^3-5p^5-p^7.
\end{align*}
Substitution of \(p=x-x^{-1}\)
reconstructs \eqref{eq:seed} exactly.

The accompanying standard-library Python
verifier uses integer Laurent polynomials
to evaluate \eqref{eq:closed}, uses
\eqref{eq:binomial} to extract \(a,b\),
and independently substitutes
\(p=x-x^{-1}\) to reconstruct each
Laurent polynomial. It obtains
\[
\begin{array}{c|r|r|r}
 L&\max\deg_x V_L&\degp V_L&b_{\degp V_L}\\ \hline
 1&8&7&-1\\
 3&12&11&-1\\
 5&16&15&-1\\
 7&20&19&-1\\
 23&52&51&-1
\end{array}
\]
and checks the exact skein identity with
the same \(H=-x-x^5\) in each case.

A separate diagram computation starts
from the five actual rational closures,
uses one uniform mirror choice to match
\eqref{eq:seed}, and evaluates their
normalized Jones polynomials with
Regina, from planar-diagram codes
generated by Spherogram.
The crossing counts are \(3,5,7,9,25\).
Every polynomial agrees coefficient
for coefficient with \eqref{eq:closed};
no additional monomial correction is
made. The diagram codes, raw outputs,
normalization and integer-arithmetic
checks accompany the note.
The general degree statement is proved
above; the finite checks provide
reproducibility, not a replacement
for that proof.

\section*{Statement on the use of AI systems}
The computation in this note was carried out within an AI collaboration directed by the author (GPT-6 through Codex and Claude Fable 5.1 through Claude Code) and verified independently by both systems; the author takes responsibility for every statement.

\begin{thebibliography}{9}
\bibitem{CC}
A.~Carmi and E.~Cohen,
\emph{On trivial Jones--Vassiliev polynomials},
arXiv:2606.22410v1, 21 June 2026.
\url{https://arxiv.org/abs/2606.22410v1}.

\bibitem{GZ}
S.~Garoufalidis and D.~Zagier,
\emph{Knots, perturbative series and quantum modularity},
SIGMA \textbf{20} (2024), 055, 87 pp.;
arXiv:2111.06645.
\href{https://doi.org/10.3842/SIGMA.2024.055}{doi:10.3842/SIGMA.2024.055}.
The initial polynomials used here appear on pp.~74 and 79.
\end{thebibliography}
\end{document}
